\documentclass[11pt,a4paper]{article}

\newcommand{\doctitle}{Documento de Visión}
\newcommand{\docsubtitle}{WARDEN - Wireless Attack Reproduction and Detection Environment}

\usepackage{warden-style}

\begin{document}

% -- Portada -------------------------------------------------------------------
\wardencoverpage{0.1}{Abril 2026}{Borrador}

% -- Historial de Revisiones ---------------------------------------------------
\begin{wardenrevisions}
\revrow{0.1}{Abril 2026}{Santiago Valencia León}{Borrador inicial}
\end{wardenrevisions}

% -- Tabla de Contenidos -------------------------------------------------------
\tableofcontents
\newpage

% ==============================================================================
\section{Introducción}
% ==============================================================================

\subsection{Propósito}
Este documento define la visión del proyecto WARDEN (Wireless Attack
Reproduction and Detection Environment), un sistema experimental dual
compuesto por un subsistema ofensivo y un subsistema defensivo, desarrollado
con fines estrictamente académicos y educativos. El proyecto reproduce, en un
ambiente de laboratorio controlado, una cadena realista de tres ataques
encadenados contra redes WiFi y, simultáneamente, los detecta mediante
captura pasiva del medio inalámbrico.

Este documento establece el contexto del problema, la solución propuesta, los
escenarios de ataque y detección, los usuarios objetivo, el alcance del
proyecto y los objetivos de alto nivel que orientan todas las decisiones de
ingeniería subsecuentes. Está dirigido a los integrantes del equipo de
desarrollo, al profesor evaluador de la asignatura, y a cualquier persona que
requiera entender qué es WARDEN y por qué existe.

\subsection{Naturaleza Académica del Proyecto}
WARDEN es un proyecto desarrollado en el contexto de la asignatura Teoría de
Sistemas de la Universidad Tecnológica de Pereira. Su finalidad es
exclusivamente pedagógica: permitir al equipo de desarrollo y a los
evaluadores del curso comprender, de manera empírica, cómo funcionan los
ataques WiFi a nivel de capa de enlace y cómo pueden ser detectados desde el
otro lado del medio inalámbrico.

El sistema se construye, prueba y ejecuta en un ambiente de laboratorio
controlado compuesto exclusivamente por equipos propiedad de los integrantes
del proyecto. En ningún momento se realizan pruebas sobre redes, dispositivos
o usuarios ajenos al equipo. El código fuente, la documentación y los
resultados experimentales tienen como destino el aprendizaje técnico y la
entrega académica de la asignatura, y no están orientados al ataque o
comprometimiento de infraestructura real.

\subsection{Alcance del Documento}
Este Documento de Visión cubre la totalidad del sistema WARDEN, incluyendo
sus dos subsistemas principales:

\begin{itemize}
  \item \textbf{Subsistema Ofensivo:} Atacante basado en ESP32 que ejecuta una
    cadena de tres ataques encadenados (Beacon Flooding, Deauthentication,
    Evil Twin con captura de credenciales).
  \item \textbf{Subsistema Defensivo:} Detector pasivo basado en el adaptador
    Panda Wireless PAU0B AC600 y software de análisis en Python, que
    identifica cada fase de la cadena ofensiva y emite alertas en tiempo real.
\end{itemize}

El detalle técnico y narrativo del subsistema ofensivo se encuentra en el
documento complementario \emph{Threat Model y Cadena de Ataque}, mientras que
los detalles de requerimientos, arquitectura y diseño están desarrollados en
los documentos correspondientes del proyecto.

\subsection{Definiciones y Acrónimos}

\begin{datatable}{L{3cm} L{12.3cm}}{Término & Definición}
\drow{WARDEN   & Wireless Attack Reproduction and Detection Environment. Nombre del proyecto.}
\drow{IDS      & Intrusion Detection System. Sistema de Detección de Intrusiones}
\drow{WIDS     & Wireless Intrusion Detection System. Variante de IDS especializada en redes inalámbricas. Categoría técnica del subsistema defensivo.}
\drow{AP       & Access Point. Punto de acceso inalámbrico}
\drow{Rogue AP & Punto de acceso no autorizado, generalmente malicioso}
\drow{SSID     & Service Set Identifier. Nombre público de una red WiFi}
\drow{BSSID    & Basic Service Set Identifier. Dirección MAC del punto de acceso}
\drow{MAC      & Media Access Control. Dirección física de un dispositivo de red}
\drow{ESP32    & Microcontrolador de bajo costo con capacidades WiFi y Bluetooth fabricado por Espressif Systems}
\drow{Modo Monitor & Modo de operación de un adaptador WiFi que permite capturar todos los frames 802.11 del entorno sin asociarse a ninguna red}
\drow{Frame 802.11 & Unidad de datos del protocolo WiFi a nivel de capa de enlace}
\drow{Evil Twin & Ataque que crea un AP falso imitando el SSID de uno legítimo}
\drow{Deauth   & Deauthentication. Frame de gestión que desconecta clientes de un AP}
\drow{Beacon   & Frame de gestión emitido periódicamente por un AP para anunciar su presencia}
\drow{Beacon Flooding & Ataque que consiste en emitir un número elevado de beacons falsos para saturar el entorno inalámbrico}
\drow{Portal Cautivo & Página web mostrada automáticamente por un cliente al conectarse a una red WiFi abierta, típicamente para autenticación o aceptación de términos}
\drow{MITRE ATT\&CK & Matriz pública de tácticas y técnicas adversarias documentadas, mantenida por MITRE Corporation}
\drow{Kill Chain & Modelo conceptual de las fases secuenciales de un ataque cibernético}
\drow{PCAP     & Packet Capture. Formato de archivo para almacenar capturas de tráfico de red}
\drow{Scapy    & Biblioteca de Python para manipulación y análisis de paquetes de red}
\drow{RF       & Radiofrecuencia. Medio físico de transmisión del protocolo WiFi}
\drow{ROE      & Rules of Engagement. Reglas de ejecución de un ejercicio ofensivo}
\end{datatable}

\subsection{Referencias}
\begin{itemize}
  \item IEEE Std 802.11-2020. Wireless LAN Medium Access Control (MAC) and
    Physical Layer (PHY) Specifications.
  \item Ley 1273 de 2009 (República de Colombia). De la protección de la
    información y de los datos.
  \item MITRE ATT\&CK Framework (\texttt{https://attack.mitre.org}).
  \item RFC 2119. Palabras clave para indicar niveles de requerimiento.
  \item Documentación oficial de Scapy (\texttt{https://scapy.net}).
  \item Documentación oficial de ESP-IDF (\texttt{https://docs.espressif.com}).
  \item Documento WARDEN \emph{Threat Model y Cadena de Ataque}.
\end{itemize}

% ==============================================================================
\section{Planteamiento del Problema}
% ==============================================================================

\subsection{Contexto}
Las redes WiFi se han convertido en el medio de conectividad predominante en
hogares, empresas e instituciones educativas. Sin embargo, la naturaleza
física del medio inalámbrico, la radiofrecuencia, expone al protocolo 802.11
a una superficie de ataque que no existe en las redes cableadas: cualquier
dispositivo con una antena dentro del rango de cobertura puede observar,
interferir o suplantar comunicaciones sin necesidad de acceso físico a la
infraestructura.

A pesar de esta exposición, la mayoría de los usuarios finales, e incluso
muchos administradores de red, desconocen los mecanismos por los cuales
ocurren los ataques a nivel de capa de enlace del protocolo 802.11. Casos
documentados como las campañas DarkHotel (Kaspersky, 2007 en adelante), los
ataques de APT28 contra hoteles de viajeros frecuentes (FireEye, 2017), y los
ejercicios de prueba de concepto presentados en conferencias como DEF CON y
Black Hat, demuestran que los ataques WiFi son una técnica activa,
contemporánea y accesible para atacantes con bajo presupuesto.

Desde la perspectiva de la teoría de sistemas, una red WiFi y su entorno
inmediato conforman un sistema abierto con múltiples componentes
interdependientes (AP, clientes, atacantes potenciales, observadores) cuya
interacción puede ser modelada, observada y caracterizada. WARDEN aborda esta
realidad construyendo dos subsistemas complementarios y simétricos: un
subsistema ofensivo que reproduce empíricamente una cadena realista de
ataques, y un subsistema defensivo que la detecta. La interacción entre
ambos, observable y mensurable en tiempo real, es el objeto de estudio del
proyecto.

\subsection{El Problema}
Los problemas específicos que aborda este proyecto son:

\begin{itemize}
  \item Los tres ataques considerados (Beacon Flooding, Deauthentication y
    Evil Twin) suelen estudiarse de forma aislada en el material académico
    disponible, perdiendo la dimensión más realista del fenómeno: en
    incidentes reales, estos ataques se ejecutan \textbf{encadenados} como
    fases de una misma operación de robo de credenciales.
  \item El material educativo sobre estos ataques tiende a centrarse
    exclusivamente en la perspectiva ofensiva (cómo se ejecuta el ataque) o
    exclusivamente en la perspectiva defensiva (cómo se detecta), pero rara
    vez integra ambas en un mismo ejercicio que permita observar la
    interacción ataque-defensa en tiempo real.
  \item La manipulación de frames 802.11 a bajo nivel con hardware común
    (por ejemplo, la tarjeta WiFi integrada de una laptop) está deliberadamente
    restringida por los drivers del sistema operativo, lo cual dificulta la
    reproducción experimental de los ataques. Esto deja un vacío didáctico
    entre los conceptos teóricos y la práctica real.
  \item No se encuentra un ejemplo académico claramente delimitado que combine
    un generador de cadena de ataques basado en un microcontrolador accesible
    (ESP32) con un detector pasivo basado en un adaptador WiFi con modo
    monitor, formando un sistema cerrado y autocontenido para fines
    educativos.
\end{itemize}

\subsection{Impacto}
La consecuencia de esta brecha es que los conceptos de seguridad inalámbrica
permanecen abstractos para estudiantes y profesionales en formación. El
desarrollo de un sistema integrado como WARDEN permite:

\begin{itemize}
  \item Materializar en hardware y software accesibles tanto la generación
    como la detección de tres ataques comunes, haciendo posible observar
    empíricamente la interacción atacante-defensor.
  \item Comprender, desde una perspectiva sistémica, cómo componentes
    heterogéneos (un microcontrolador ofensivo, dispositivos víctima, un
    adaptador WiFi especializado y un computador defensor) pueden integrarse
    para modelar un escenario de ciberseguridad completo.
  \item Estudiar la lógica encadenada de un ataque real, donde las fases
    individuales toman significado solo cuando se observan en su secuencia
    completa.
  \item Generar un recurso didáctico reproducible que puede ser replicado,
    extendido o utilizado como base para trabajos futuros en seguridad
    inalámbrica.
\end{itemize}

% ==============================================================================
\section{Declaración de Visión}
% ==============================================================================

Para estudiantes y profesionales en formación interesados en comprender la
seguridad de redes inalámbricas desde una perspectiva empírica e integral,
\textbf{WARDEN} es un sistema experimental dual de ataque y detección que
permite, mediante hardware accesible (una ESP32 como generador de la cadena
de ataques y un adaptador Panda Wireless PAU0B AC600 como sensor pasivo),
reproducir y detectar una cadena realista de tres ataques encadenados contra
redes WiFi: Beacon Flooding, Deauthentication y Evil Twin con robo de
credenciales mediante portal cautivo falso.

A diferencia de suites de pentesting profesional orientadas únicamente al
ataque (como Aircrack-ng, Wifiphisher o Kismet en sus modos ofensivos),
WARDEN integra simétricamente la perspectiva defensiva, permitiendo observar
en tiempo real cómo cada fase del ataque produce evidencias detectables. A
diferencia de productos WIDS comerciales (Cisco CleanAir, Aruba RFProtect,
Fortinet FortiGate Wireless), WARDEN no busca desempeño empresarial ni
cobertura exhaustiva, sino servir como un modelo pedagógico claro,
reproducible y de bajo costo donde el ataque y la detección son ambos
ciudadanos de primera clase.

WARDEN opera bajo dos principios fundamentales: \textbf{contención total
del laboratorio} (todos los ataques se dirigen exclusivamente a equipos
propiedad del equipo) y \textbf{detección estrictamente pasiva} (el
detector solo escucha el medio, jamás inyecta tráfico). Las alertas y
resultados se emiten por consola con la información mínima necesaria para
identificar cada fase del ataque y sus metadatos asociados.

% ==============================================================================
\section{Subsistemas Principales}
% ==============================================================================

WARDEN se compone de dos subsistemas funcionalmente independientes pero
diseñados para operar de forma complementaria.

\subsection{Subsistema Ofensivo}

\textbf{Componente principal:} Microcontrolador ESP32 con firmware
desarrollado por el equipo.

\textbf{Función:} Reproducir una cadena de ataques WiFi en tres fases,
dirigida exclusivamente a los dispositivos víctima del laboratorio.

\textbf{Cadena de ataques (resumen):}

\begin{enumerate}
  \item \textbf{Fase 1 - Beacon Flooding (reconocimiento y distracción):} La
    ESP32 emite un número elevado de beacons falsos con SSIDs aleatorios para
    saturar el entorno inalámbrico de la víctima.
  \item \textbf{Fase 2 - Deauthentication (aislamiento de la víctima):} La
    ESP32 envía frames de deautenticación dirigidos al BSSID del router
    legítimo de laboratorio, forzando la desconexión de los clientes víctima
    (Redmi Note 7 y Dell Latitude E6220).
  \item \textbf{Fase 3 - Evil Twin con portal cautivo falso (explotación):}
    La ESP32 crea un AP abierto que imita el SSID de la red legítima. Cuando
    la víctima se reconecta automáticamente, se le presenta un portal cautivo
    falso desarrollado por el equipo, que captura las credenciales ingresadas
    (con datos ficticios para fines exclusivamente demostrativos).
\end{enumerate}

\textbf{Atacante simulado:} Actor oportunista cercano operando en un espacio
público compartido, con motivación de robo de credenciales mediante
suplantación de AP. Este perfil corresponde a campañas reales documentadas
como DarkHotel y APT28.

\textbf{Mapeo a MITRE ATT\&CK:}

\begin{datatable}{L{2.5cm} L{4cm} L{8.8cm}}{Técnica & Nombre & Aplicación en WARDEN}
\drow{T1498 & Network Denial of Service & Fase 1: Beacon Flooding satura el entorno}
\drow{T1499 & Endpoint Denial of Service & Fase 2: Deauth desconecta clientes legítimos}
\drow{T1557 & Adversary-in-the-Middle    & Fase 3: Evil Twin se posiciona entre cliente y red}
\drow{T1056.003 & Web Portal Capture     & Fase 3: Captura de credenciales por portal cautivo}
\end{datatable}

El detalle completo del modelado de amenaza, narrativa del atacante, técnicas
empleadas y reglas de ejecución se encuentra en el documento WARDEN
\emph{Threat Model y Cadena de Ataque}.

\subsection{Subsistema Defensivo}

\textbf{Componente principal:} Estación de detección compuesta por una laptop
con Arch Linux y el adaptador Panda Wireless PAU0B AC600 en modo monitor.

\textbf{Función:} Capturar pasivamente el espectro 802.11 en la banda 2.4 GHz
e identificar la firma de cada fase de la cadena ofensiva en tiempo real,
emitiendo alertas estructuradas por consola.

\textbf{Capacidades de detección:}

\begin{enumerate}
  \item \textbf{Detección de Beacon Flooding:} Identificación de tasas
    anómalas de beacons únicos por segundo y detección de rangos sospechosos
    de BSSIDs falsos.
  \item \textbf{Detección de Deauthentication:} Identificación de frames de
    deauth dirigidos al BSSID protegido, detección de tasas anómalas y de
    direccionalidad sospechosa.
  \item \textbf{Detección de Evil Twin:} Identificación de SSIDs duplicados
    con BSSIDs distintos al BSSID legítimo conocido, detección de cambios
    repentinos en las características de seguridad anunciadas (por ejemplo,
    SSID que normalmente requiere WPA2 apareciendo como abierto).
  \item \textbf{Correlación de cadena (objetivo deseable):} Detección de
    patrones que sugieren ataques encadenados: Beacon Flood seguido de Deauth
    seguido de aparición de Evil Twin en un intervalo corto.
\end{enumerate}

\textbf{Categoría técnica:} WIDS pasivo basado en firmas y heurísticas
explícitas. El subsistema defensivo no es una respuesta automatizada (no es
un IPS); su función se limita a detectar y reportar.

% ==============================================================================
\section{Escenario del Laboratorio}
% ==============================================================================

\subsection{Topología Física}

El laboratorio WARDEN está compuesto por cinco elementos físicos
interactuando exclusivamente a través del medio inalámbrico en la banda 2.4
GHz.

\begin{datatable}{L{4cm} L{4.5cm} L{6.8cm}}{Elemento & Rol & Descripción}
\drow{Router de laboratorio & Red legítima objetivo & Router económico adquirido específicamente para el proyecto, configurado con SSID \texttt{LAB\_WARDEN\_UTP} y desconectado de internet. BSSID protegido por el subsistema defensivo.}
\drow{Redmi Note 7 (MIUI 12.5) & Víctima primaria & Cliente Android conectado al router legítimo. Recibe la cadena completa de ataques. Su comportamiento de reconexión automática y manejo de portales cautivos lo hacen el objetivo más realista del escenario.}
\drow{Dell Latitude E6220 (Ubuntu Server 22.04.5 LTS) & Víctima secundaria / testigo & Cliente Linux conectado al router legítimo. Recibe Beacon Flood y Deauth. Sus logs del sistema (\texttt{dmesg}, \texttt{journalctl}) son evidencia técnica de los ataques sufridos.}
\drow{ESP32 & Atacante & Microcontrolador ejecutando el firmware ofensivo de WARDEN. Genera la cadena de ataques contra los BSSID y MAC del laboratorio.}
\drow{Laptop Arch Linux + Panda PAU0B & Detector & Estación con adaptador en modo monitor capturando pasivamente todo el espectro 802.11. Ejecuta el software detector en Python con Scapy.}
\end{datatable}

\subsection{Diagrama Conceptual}

La interacción entre los elementos puede representarse así:

\vspace{0.3cm}
\begin{verbatim}
+-------------------+         +-----------------------+
|  ROUTER LEGITIMO  | <-----> |  REDMI NOTE 7         |
|  LAB_WARDEN_UTP   |   WiFi  |  (victima primaria)   |
|  BSSID protegido  |         +-----------------------+
+-------------------+
         ^                    +-----------------------+
         |                    |  DELL LATITUDE E6220  |
         |             <----> |  (victima secundaria) |
         |                    +-----------------------+
         |
         | (Cadena ofensiva por radiofrecuencia)
         |
+--------+----------+
|      ESP32        |
|     ATACANTE      |
|  Fase 1: Beacons  |
|  Fase 2: Deauth   |
|  Fase 3: Evil Twin|
+-------------------+

       +================================================+
       |  ARCH LINUX + PANDA PAU0B (modo monitor)       |
       |  DETECTOR PASIVO                                |
       |  Captura todo el espectro, no inyecta nada     |
       |  Emite alertas por fase y de la cadena completa|
       +================================================+
\end{verbatim}

\subsection{Banda y Canal}

Toda la operación se realiza en la banda de 2.4 GHz, en un canal específico
seleccionado durante la configuración del router de laboratorio. La elección
del canal busca minimizar la coincidencia con redes WiFi vecinas.

% ==============================================================================
\section{Interesados y Usuarios Objetivo}
% ==============================================================================

\subsection{Interesados (Stakeholders)}

\begin{datatable}{L{4cm} L{3.5cm} L{8.0cm}}{Interesado & Rol & Interés}
\drow{Santiago Valencia León  & Integrante del equipo  & Desarrollo del sistema, aprendizaje técnico en seguridad WiFi y teoría de sistemas}
\drow{[Integrante 2]          & Integrante del equipo  & Desarrollo del sistema, aprendizaje técnico en seguridad WiFi y teoría de sistemas}
\drow{[Integrante 3]          & Integrante del equipo  & Desarrollo del sistema, aprendizaje técnico en seguridad WiFi y teoría de sistemas}
\drow{Juan Andrés García Moreno & Profesor evaluador   & Evaluación del proyecto como cumplimiento de los objetivos de la asignatura Teoría de Sistemas}
\drow{UTP, Facultad de Ingeniería & Institución académica & Formación de estudiantes en el análisis sistémico de problemas de ingeniería}
\end{datatable}

\subsection{Usuarios Objetivo}

\begin{datatable}{L{3.5cm} L{5.5cm} L{6.3cm}}{Rol de Usuario & Descripción & Interacción Principal}
\drow{Operador del Atacante & Persona que controla la ESP32 y selecciona qué fase de la cadena ejecutar. Inicia la cadena completa o cada fase individualmente. & Botón físico, selector o consola serial de la ESP32}
\drow{Operador del Detector & Persona que ejecuta el software detector en la laptop con el Panda. Monitorea las alertas emitidas por consola en tiempo real. & Línea de comandos de la laptop detectora}
\drow{Operador de la Víctima & Persona que opera los dispositivos víctima durante la demostración. Para fines de demostración, ingresa credenciales ficticias en el portal cautivo falso. & Interfaz nativa del Redmi Note 7 y de la Dell Latitude}
\drow{Observador & Persona que estudia el funcionamiento del sistema sin interactuar directamente. Es el rol principal de los estudiantes espectadores y del profesor evaluador durante la demostración. & Observación del montaje completo y lectura de las salidas del detector y del atacante}
\end{datatable}

% ==============================================================================
\section{Alcance del Producto}
% ==============================================================================

\subsection{Dentro del Alcance}
\begin{itemize}
  \item \textbf{Subsistema Ofensivo (ESP32):} Firmware capaz de ejecutar la
    cadena completa de tres ataques (Beacon Flooding, Deauthentication y Evil
    Twin con portal cautivo falso) tanto de forma encadenada automática como
    fase por fase manual. La cadena se dirige exclusivamente al BSSID y MAC
    del laboratorio.
  \item \textbf{Portal Cautivo Falso:} Página web servida por la ESP32 durante
    la fase de Evil Twin, que solicita credenciales ficticias bajo pretexto
    de re-autenticación. Las credenciales ingresadas se registran por consola
    serial como evidencia de captura exitosa, con fines exclusivamente
    demostrativos.
  \item \textbf{Infraestructura de Laboratorio:} Red WiFi dedicada compuesta
    por un router económico desconectado de internet con SSID
    \texttt{LAB\_WARDEN\_UTP} y dos dispositivos víctima controlados por el
    equipo (Redmi Note 7 con MIUI 12.5 y Dell Latitude E6220 con Ubuntu Server
    22.04.5 LTS).
  \item \textbf{Subsistema Defensivo:} Software en Python basado en Scapy que
    captura frames 802.11 en tiempo real desde el adaptador Panda Wireless
    PAU0B AC600 en modo monitor y detecta cada una de las tres fases de la
    cadena ofensiva.
  \item \textbf{Correlación de Cadena (objetivo deseable):} Lógica de
    correlación temporal en el detector que identifica las tres fases como
    una secuencia coordinada y emite una alerta de cadena completa.
  \item \textbf{Alertas por Consola:} Emisión de alertas en formato legible
    por humanos al detectarse cada fase, incluyendo marca de tiempo, tipo de
    ataque, BSSID afectado, MAC de origen, y contadores cuando apliquen.
  \item \textbf{Operación en 2.4 GHz:} Todas las pruebas se realizan en la
    banda de 2.4 GHz para garantizar compatibilidad con el chipset MT7610U.
  \item \textbf{Capturas PCAP de Referencia:} Generación de un conjunto de
    archivos PCAP con cadenas de ataque representativas, utilizados tanto
    para testing del detector como para desarrollo por parte de los
    integrantes que no disponen del adaptador Panda.
  \item \textbf{Documentación de Ingeniería:} Documento de Visión, Threat
    Model y Cadena de Ataque, Marco Teórico, Especificación de Requerimientos,
    Documento de Casos de Uso, Documento de Arquitectura y Documento de
    Diseño Detallado.
  \item \textbf{Demostración Funcional:} Demostración en vivo del sistema
    completo end-to-end durante la entrega del proyecto, incluyendo la
    cadena de ataque y la detección simultánea.
\end{itemize}

\subsection{Fuera del Alcance}
\begin{itemize}
  \item Detección o ejecución de ataques WiFi distintos a los tres
    encadenados (por ejemplo: KARMA, ataques contra WPS, captura de
    handshakes WPA, ataques contra PMKID, downgrade de cifrado).
  \item Captura real o uso real de credenciales. El portal cautivo solo
    captura datos ficticios introducidos durante demostraciones controladas.
  \item Crackeo offline de contraseñas o ataques de fuerza bruta sobre
    handshakes capturados.
  \item Interfaz gráfica de usuario en cualquiera de los subsistemas (la
    operación es exclusivamente por línea de comandos y, en el caso del
    portal cautivo, por una página web mínima funcional).
  \item Prevención, bloqueo o mitigación automática de los ataques
    detectados. El subsistema defensivo es estrictamente de detección y
    reporte.
  \item Captura y análisis en la banda de 5 GHz.
  \item Soporte para redes WPA3.
  \item Técnicas de detección basadas en aprendizaje automático o
    estadísticas avanzadas. Solo se utilizan reglas heurísticas explícitas.
  \item Despliegue en entornos de producción.
  \item Integración con sistemas SIEM u otras plataformas de gestión de
    eventos de seguridad.
  \item Análisis forense posterior a los ataques.
  \item Captura distribuida mediante múltiples sensores.
  \item Soporte para sistemas operativos distintos a Linux en el subsistema
    defensivo.
\end{itemize}

% ==============================================================================
\section{Objetivos y Criterios de Éxito}
% ==============================================================================

\subsection{Objetivos del Subsistema Ofensivo}

\begin{datatable}{L{1.4cm} L{10.3cm} L{1.9cm}}{ID & Objetivo & Prioridad}
\drow{O-01 & Implementar Beacon Flooding capaz de emitir al menos 50 beacons únicos por segundo durante la fase 1 & Alta}
\drow{O-02 & Implementar Deauthentication dirigido al BSSID del router de laboratorio, capaz de desconectar a las víctimas en menos de 5 segundos & Alta}
\drow{O-03 & Implementar Evil Twin que clone el SSID del router de laboratorio y se anuncie como red abierta & Alta}
\drow{O-04 & Implementar portal cautivo falso que registre credenciales ficticias ingresadas durante la fase 3 & Alta}
\drow{O-05 & Permitir ejecución encadenada automática y ejecución manual fase por fase & Alta}
\drow{O-06 & Mapear cada fase a una técnica documentada de la matriz MITRE ATT\&CK & Media}
\end{datatable}

\subsection{Objetivos del Subsistema Defensivo}

\begin{datatable}{L{1.4cm} L{10.3cm} L{1.9cm}}{ID & Objetivo & Prioridad}
\drow{D-01 & Detectar la fase de Beacon Flooding con precisión superior al 90\% en el laboratorio & Alta}
\drow{D-02 & Detectar la fase de Deauthentication con precisión superior al 90\% en el laboratorio & Alta}
\drow{D-03 & Detectar la fase de Evil Twin identificando SSIDs duplicados con BSSID distinto al protegido & Alta}
\drow{D-04 & Emitir alertas en consola con marca de tiempo, tipo de ataque y metadatos relevantes en menos de 2 segundos desde la detección & Alta}
\drow{D-05 & Operar en modo completamente pasivo, sin inyectar tráfico al medio inalámbrico & Alta}
\drow{D-06 & Soportar análisis tanto en captura en vivo como sobre archivos PCAP para fines de desarrollo y testing & Alta}
\drow{D-07 & Identificar correlación temporal entre las tres fases y emitir alerta de cadena completa cuando ocurran en secuencia & Media}
\end{datatable}

\subsection{Objetivos Generales del Proyecto}

\begin{datatable}{L{1.4cm} L{10.3cm} L{1.9cm}}{ID & Objetivo & Prioridad}
\drow{G-01 & Producir la documentación completa de ingeniería del proyecto & Alta}
\drow{G-02 & Realizar una demostración funcional end-to-end durante la entrega, mostrando la cadena ofensiva y la detección simultánea & Alta}
\drow{G-03 & Generar capturas PCAP de referencia para evaluación y reproducibilidad & Media}
\drow{G-04 & Mantener el código fuente y la documentación en un repositorio Git público & Media}
\end{datatable}

\subsection{Criterios de Éxito}
\begin{itemize}
  \item Los objetivos O-01 a O-05, D-01 a D-06, y G-01 a G-02 se cumplen y
    se demuestran en vivo durante la sesión de entrega.
  \item Durante la demostración, la cadena de ataque ejecutada por la ESP32
    es detectada por el subsistema defensivo en al menos 9 de 10 ejecuciones.
  \item El portal cautivo falso captura credenciales ficticias durante la
    demostración y las muestra por consola serial.
  \item Los seis documentos de ingeniería están completos, revisados y
    compilados en formato PDF al momento de la entrega.
  \item La documentación y el código permiten que cualquier integrante del
    equipo responda preguntas técnicas fundamentales del sistema sin
    necesidad de consultar a otro integrante.
\end{itemize}

% ==============================================================================
\section{Restricciones}
% ==============================================================================

\subsection{Restricciones Técnicas}
\begin{itemize}
  \item El sistema operativo del subsistema defensivo debe ser Linux. Se han
    validado dos distribuciones: Arch Linux (estación principal) y Ubuntu
    Server 22.04.5 LTS (dispositivo víctima secundario).
  \item El adaptador WiFi para captura debe soportar modo monitor. Se utiliza
    el Panda Wireless PAU0B AC600 con chipset MediaTek MT7610U.
  \item La operación se limita a la banda de 2.4 GHz debido a limitaciones
    conocidas de estabilidad del driver MT7610U en 5 GHz para modo monitor.
  \item El subsistema ofensivo utiliza una ESP32 estándar. No se requiere una
    variante específica (ESP32-S3, ESP32-C3, etc.).
  \item El lenguaje de programación del subsistema defensivo es Python 3
    utilizando la biblioteca Scapy.
  \item El firmware de la ESP32 se desarrolla utilizando Arduino-ESP32 o
    ESP-IDF.
  \item El portal cautivo falso se sirve desde la ESP32 mediante un servidor
    web embebido y un servidor DNS local que redirige todas las consultas a
    la ESP32.
\end{itemize}

\subsection{Restricciones de Diseño}
\begin{itemize}
  \item El subsistema defensivo debe ser completamente pasivo: no puede
    inyectar frames ni modificar de ninguna forma el tráfico del medio
    inalámbrico.
  \item La detección debe realizarse en tiempo real sobre el flujo de
    captura, no sobre archivos PCAP almacenados. No obstante, el mismo código
    del detector debe poder operar sobre archivos PCAP para fines de
    desarrollo y testing.
  \item Los ataques generados por la ESP32 deben dirigirse exclusivamente al
    BSSID del router de laboratorio o a las MAC de los dispositivos víctima
    del equipo. Bajo ninguna circunstancia se dirigirán a redes o
    dispositivos ajenos al proyecto.
  \item El portal cautivo falso debe ser claramente identificado como un
    componente de demostración educativa, y solo aceptará datos ficticios
    durante las pruebas.
  \item El sistema completo debe funcionar sin conexión a internet.
\end{itemize}

\subsection{Restricciones de Recursos}
\begin{itemize}
  \item El proyecto debe completarse en un plazo de tres a cuatro semanas a
    partir de la semana 12 del semestre académico.
  \item El equipo está conformado por tres integrantes.
  \item El presupuesto para adquisición de hardware adicional (router de
    laboratorio) es asumido por el equipo.
  \item Solo un integrante del equipo dispone del adaptador Panda Wireless
    PAU0B AC600. En consecuencia, las capturas en vivo solo pueden realizarse
    desde su estación.
\end{itemize}

\subsection{Restricciones Éticas y Legales}
El presente proyecto involucra la ejecución de ataques contra redes WiFi en
un entorno experimental con finalidad estrictamente educativa. Dado que la
interceptación de comunicaciones inalámbricas y el acceso no autorizado a
sistemas informáticos están tipificados como delitos en la legislación
colombiana bajo la Ley 1273 de 2009 (artículos 269A sobre acceso abusivo a
sistema informático, 269C sobre interceptación de datos informáticos y 269D
sobre daño informático), resulta indispensable ejecutar la totalidad de las
pruebas en un ambiente controlado que garantice que ninguna red, dispositivo
o usuario ajeno al proyecto se vea afectado.

\textbf{Definición del ambiente controlado.} A diferencia de otros escenarios
de ciberseguridad donde el aislamiento se logra mediante máquinas virtuales,
en este proyecto el aislamiento mediante virtualización no es aplicable. Esto
se debe a que los ataques objeto de estudio operan sobre la capa de enlace
del protocolo 802.11, es decir, sobre frames de gestión y control que se
transmiten exclusivamente a través del medio físico de radiofrecuencia. Las
máquinas virtuales no poseen interfaces de radio reales, solo interfaces de
red virtualizadas a nivel de IP, por lo que no pueden emitir ni capturar
frames 802.11 reales.

En consecuencia, el ambiente controlado para este proyecto se define
mediante \textbf{aislamiento físico y lógico}, compuesto por los siguientes
elementos:

\begin{enumerate}
  \item \textbf{Red WiFi dedicada al laboratorio.} Se utiliza un router WiFi
    adquirido específicamente para el proyecto, configurado con un SSID
    identificable como entorno de prueba (\texttt{LAB\_WARDEN\_UTP}) y
    desconectado de cualquier enlace a internet. Esto garantiza que ningún
    tráfico externo interfiera con las capturas y que ningún dispositivo
    ajeno se conecte accidentalmente.
  \item \textbf{Dispositivos víctima de propiedad del equipo.} Todos los
    dispositivos que actúan como clientes legítimos de la red atacada son
    propiedad de los integrantes del proyecto: un celular Redmi Note 7 con
    MIUI 12.5 y una laptop Dell Latitude E6220 con Ubuntu Server 22.04.5
    LTS. No se realiza ninguna prueba contra dispositivos de terceros.
  \item \textbf{Dispositivo atacante bajo control directo.} La ESP32 ejecuta
    firmware desarrollado por el equipo, cuyos ataques están dirigidos
    exclusivamente al BSSID del router de laboratorio y a las MAC de las
    víctimas del equipo. Los frames maliciosos no se dirigen a ninguna otra
    red presente en el entorno.
  \item \textbf{Direccionalidad de los ataques.} Las tres fases de la cadena
    se configuran para afectar únicamente al BSSID específico del router de
    laboratorio y al rango de direcciones MAC de los dispositivos víctima
    del equipo. Aunque los frames se emiten por radiofrecuencia y son
    físicamente audibles por otros receptores cercanos, no generan efectos
    sobre redes o dispositivos ajenos, pues no están dirigidos a ellos.
  \item \textbf{Datos ficticios en el portal cautivo.} El portal cautivo
    falso solo recibe datos ficticios introducidos por miembros del equipo
    durante demostraciones controladas. Bajo ninguna circunstancia se
    introducen credenciales reales.
  \item \textbf{Ubicación de las pruebas.} Las sesiones de laboratorio se
    realizan en espacios donde la densidad de redes WiFi vecinas sea mínima,
    como medida de precaución adicional.
  \item \textbf{Registro de sesiones.} Cada sesión de pruebas se documenta
    con fecha, hora de inicio y fin, participantes, fase de la cadena
    ejecutada y resultados obtenidos.
\end{enumerate}

% ==============================================================================
\section{Suposiciones y Dependencias}
% ==============================================================================

\subsection{Suposiciones}
\begin{itemize}
  \item El chipset MediaTek MT7610U del adaptador Panda Wireless PAU0B AC600
    mantiene compatibilidad con modo monitor en la banda 2.4 GHz bajo las
    distribuciones Linux utilizadas.
  \item La ESP32 disponible permite, mediante bibliotecas existentes o
    modificaciones del SDK, la emisión de frames 802.11 de gestión
    arbitrarios (incluidos frames de deauth y beacons falsos), así como
    operar simultáneamente como AP para servir el portal cautivo.
  \item Los dispositivos víctima (Redmi Note 7 y Dell Latitude E6220)
    reaccionan de forma observable a los ataques: se desconectan ante el
    deauth y, en el caso del celular, intentan reconexión automática y
    presentan el portal cautivo al detectar una red abierta.
  \item El router de laboratorio a adquirir soporta WiFi 2.4 GHz y permite
    configurar SSID y contraseña de forma manual.
\end{itemize}

\subsection{Dependencias Externas}

\begin{datatable}{L{3cm} L{1.8cm} L{5.5cm} L{5cm}}{Dependencia & Versión & Propósito & Componente}
\drow{Python         & 3.10+   & Lenguaje del subsistema defensivo   & Detector}
\drow{Scapy          & 2.5+    & Captura y análisis de frames 802.11 & Detector}
\drow{aircrack-ng    & 1.7+    & Activación de modo monitor (airmon-ng) & Detector}
\drow{Arduino-ESP32  & 2.0+    & Framework de desarrollo del firmware & Atacante}
\drow{ESP-IDF        & 5.0+    & Alternativa al framework Arduino    & Atacante}
\drow{Linux Kernel   & 5.15+   & Soporte de drivers MT7610U          & Detector}
\drow{Git / GitHub   & ---     & Control de versiones y colaboración & Todo el proyecto}
\drow{LaTeX          & ---     & Compilación de la documentación     & Documentación}
\end{datatable}

% ==============================================================================
\section{Consideraciones Futuras}
% ==============================================================================

Los siguientes elementos están intencionalmente fuera del alcance de la
versión inicial pero se reconocen como posibles direcciones futuras de
extensión del proyecto:

\begin{itemize}
  \item Expansión de la cadena ofensiva con nuevos ataques (KARMA, captura
    de handshakes WPA, ataques contra WPS, ataques de downgrade de cifrado).
  \item Sustitución de las reglas heurísticas del detector por modelos de
    aprendizaje automático entrenados sobre capturas reales.
  \item Implementación de una interfaz gráfica o dashboard web para
    visualización en tiempo real tanto del lado ofensivo como del defensivo.
  \item Soporte para captura y análisis en la banda de 5 GHz.
  \item Integración del detector con sistemas SIEM mediante formatos
    estándar (CEF, Syslog).
  \item Extensión a sensores distribuidos coordinados para cobertura de
    áreas más amplias.
  \item Evolución del subsistema defensivo a un WIPS (Wireless Intrusion
    Prevention System) con capacidad de respuesta automática ante ataques
    detectados.
  \item Generación automática de reportes forenses a partir de las alertas
    acumuladas durante una sesión.
  \item Variantes adicionales del portal cautivo que reproduzcan estilos
    visuales de proveedores reales de WiFi público.
\end{itemize}

\end{document}
